\section{Related Work}
\label{sec:related}

\subsection{Few-Shot and Cross-Domain Detection}

Few-shot object detection learns novel categories from scarce annotations. Early meta-learning approaches condition a detector on support features, whereas transfer-learning methods show that carefully constrained fine-tuning can be surprisingly competitive. Representative methods include TFA~\cite{tfa}, FSCE~\cite{fsce}, and DeFRCN~\cite{defrcn}, which study restricted classifier adaptation, contrastive proposal representations, and decoupled classification and localization, respectively. These methods generally assume that base and novel data originate from the same domain.

CD-FSOD further introduces a severe source--target domain shift. CD-ViTO~\cite{cdvito} established a six-domain benchmark and enhanced an open-set detector with target instances. Subsequent work caches target features~\cite{ifc}, searches for foundation-model augmentations~\cite{ets}, retrieves compositional data~\cite{domainrag}, or exploits style augmentation and sparse target instances~\cite{styleproto,sivito}. VFM-MoE~\cite{vfmmoe} combines DETR with foundation-model experts, while Yu \etal~\cite{yu2026} show that progressive schedules and decoder diversity are critical to stable few-shot fine-tuning. Target-Domain Astigmatism~\cite{astigmatism}, in contrast, diagnoses dispersed attention and refines both foreground and peripheral context. Our work is complementary to these approaches: we examine whether each support-derived factor receives decision authority commensurate with its reliability.

\subsection{Vision-Language Detectors and Support Knowledge}

Open-vocabulary detectors transfer category semantics by aligning visual regions with language. Representative directions include region-text knowledge distillation in ViLD~\cite{vild}, region-level language-image pretraining in RegionCLIP~\cite{regionclip}, end-to-end multimodal decoding in MDETR~\cite{mdetr}, conditional matching in OV-DETR~\cite{ovdetr}, and image-text prompting in OWL-ViT~\cite{owlvit}. On the detection side, Deformable DETR~\cite{deformable_detr}, DAB-DETR~\cite{dab_detr}, and DN-DETR~\cite{dn_detr} improve the convergence and query design of the original DETR formulation~\cite{detr}. GLIP~\cite{glip} formulates detection as phrase grounding, while Grounding DINO~\cite{groundingdino} combines encoder-level feature fusion, language-guided query selection, and cross-modal decoder interaction with the DINO detector~\cite{dino}. Such models provide strong initializations for CD-FSOD because category names remain meaningful across domains; however, text alone cannot capture target-specific appearance.

Recent methods therefore inject visual support knowledge. LMP~\cite{lmp} constructs class prototypes and query-derived hard negatives in a visual branch, while PET-DINO~\cite{petdino} unifies instance- and category-level visual cues through prompt-enriched training. Target-Domain Astigmatism~\cite{astigmatism} jointly models foreground prototypes, contextual background, and text. These approaches demonstrate that support features can complement language, but they typically optimize each prototype as a complete representation. In contrast, CHSD distinguishes the component that defines category identity from the component that captures necessary domain appearance. ATAR further controls the stage at which, and the magnitude with which, each component enters the detector.

\subsection{Context Sensitivity and Prediction Reliability}

Context can support object-boundary estimation and category disambiguation, yet it can also become a shortcut that produces background false positives under distribution shift. Existing CD-FSOD methods use negative regions as contrastive cues~\cite{lmp}, explicitly model peripheral background~\cite{astigmatism}, or evaluate robustness on mixed-domain queries~\cite{yu2026}. Confidence calibration offers a complementary perspective: a detector should not only be accurate, but also align confidence with correctness and identify high-risk predictions~\cite{calibration,selective}.

Our focus differs in two respects. First, we intervene on the support set rather than the query: the query remains fixed while the context outside the protected support foreground changes. Second, the intervention serves as a diagnostic probe rather than as additional labeled data. FFCP measures prediction drift across support conditions, CHSD excludes context-sensitive directions from target-appearance knowledge, and ATAR contracts Soft classification authority when support evidence is unstable. Together, these modules connect support-side sensitivity to detection accuracy, calibration, and selective risk.
